\documentclass[12pt,a4paper]{article}
\usepackage[utf8]{inputenc}
\usepackage[vietnamese]{babel}
\usepackage[left=2cm,right=2cm,top=2cm,bottom=2cm]{geometry}
\usepackage{mathptmx} % Font Times New Roman
\usepackage{enumitem}

\pagestyle{plain}
\linespread{1.15}

\begin{document}

% --- HEADER QUỐC HIỆU & CƠ QUAN BAN HÀNH ---
\noindent
\begin{minipage}[t]{0.45\textwidth}
    \begin{center}
        \small
        CÔNG AN THÀNH PHỐ HÀ NỘI\\
        \textbf{\underline{PHÒNG CẢNH SÁT HÌNH SỰ}}\\[0.1cm]
        Số: 02/BB-KNHT\\
        \textit{V/v khám nghiệm hiện trường vụ án}
    \end{center}
\end{minipage}
\hfill
\begin{minipage}[t]{0.52\textwidth}
    \begin{center}
        \small
        \textbf{CỘNG HÒA XÃ HỘI CHỦ NGHĨA VIỆT NAM}\\
        \textbf{\underline{Độc lập - Tự do - Hạnh phúc}}\\[0.1cm]
        \textit{Hà Nội, ngày 25 tháng 07 năm 2026}
    \end{center}
\end{minipage}

\vspace{0.6cm}

% --- TIÊU ĐỀ VĂN BẢN ---
\begin{center}
    \textbf{\large BIÊN BẢN KHÁM NGHIỆM HIỆN TRƯỜNG}\\
    \textit{(Vụ án mạng tại số 14 Đường Bờ Sông, Phân khu Cảng)}
\end{center}

\vspace{0.3cm}

\noindent Căn cứ Điều 201 Bộ luật Tố tụng hình sự năm 2015;\\
Vào hồi 08 giờ 30 phút, ngày 25 tháng 07 năm 2026, tại căn nhà số 14, Đường Bờ Sông, Phường Phân khu Cảng, TP. Hà Nội.\\[0.1cm]
Chúng tôi gồm:
\begin{enumerate}[label=\arabic*., leftmargin=1.5cm]
    \item Đơn vị chủ trì: Đại úy Lê Minh — Điều tra viên, Phòng Cảnh sát Hình sự.
    \item Đội ngũ kỹ thuật hình sự: Trung úy Phạm Thanh — Kỹ thuật viên Phòng KTHS.
    \item Đại diện Viện Kiểm sát: Ông Nguyễn Văn Hòa — Kiểm sát viên VKSND TP. Hà Nội.
    \item Người chứng kiến: Bà Trần Thị Mai (sinh năm 1970) — Tổ trưởng dân phố số 4.
\end{enumerate}

\noindent Tiến hành khám nghiệm hiện trường vụ án hình sự theo quy định. Kết quả ghi nhận như sau:

\section*{I. MÔ TẢ TỔNG QUAN HIỆN TRƯỜNG}
Căn nhà 1 tầng cũ diện tích khoảng 75m2, nằm trong khu vực giải tỏa đền bù hẻo lánh. Căn nhà không trang bị hệ thống camera an ninh.\\
Cửa chính khóa xích ngoắc lỏng; cửa sau (đi ra hẻm nhỏ hướng ra sông) bị bật chốt trong, có dấu vết bị đẩy mở từ bên ngoài.

\section*{II. DẤU VẾT TẠI PHÒNG KHÁCH (NƠI PHÁT HIỆN THI THỂ)}
\begin{enumerate}[label=\arabic*., leftmargin=1.5cm]
    \item \textbf{Vị trí thi thể \& Vũng máu:} Nạn nhân Nguyễn Văn Khang nằm ngửa trên sàn gạch cũ cạnh bàn trà gỗ. Vũng máu đầm đìa lan từ vùng cổ nạn nhân ra xung quanh.
    \item \textbf{Dấu vết xáo trộn \& Vật chứng rơi vỡ:} Bộ bình cốc thủy tinh uống trà bị đập vỡ vụn dưới sàn nhà gần mép bàn trà. Cạnh đó phát hiện 01 hũ thủy tinh đựng hoa cúc sấy khô dán note viết tay của Hà.
    \item \textbf{Dấu vết vật chứng dính máu:} Tìm thấy 01 mảnh thủy tinh vỡ sắc nhọn dài khoảng 8cm dính máu khô và có dấu ấn vân tay dạng miết sát cạnh thi thể (ký hiệu vật chứng \texttt{p3}).
    \item \textbf{Cúc áo đứt dưới gầm bàn:} Thu giữ 01 chiếc cúc áo sơ mi nam bằng nhựa màu xanh đen bị giật đứt chỉ rơi dưới gầm bàn trà phòng khách (ký hiệu vật chứng \texttt{EV-SHIRT-BUTTON}).
    \item \textbf{Vật chứng kim loại dưới chân bàn:} Thu giữ 01 chiếc \textbf{Ghim cài áo (Huy hiệu nhận diện công ty xây dựng)} bằng đồng dính bụi vôi vữa rơi sát chân bàn trà (ký hiệu vật chứng \texttt{p8}).
    \item \textbf{Mảnh báo cũ bị xé rách:} Nhiều mảnh giấy báo cũ ố vàng bị xé rách vụn nằm rải rác trên sàn gần bàn trà (ký hiệu vật chứng \texttt{p5}).
    \item \textbf{Sổ nợ rách \& Giấy tờ tranh chấp văng vãi:} Trên bàn làm việc phát hiện 01 cuốn sổ da màu đen ghi chép cho vay nợ bị xé toạc mất 3 trang giữa. Dưới sàn nhà văng vãi tập hồ sơ gồm: \textit{Đơn tố cáo lừa đảo chiếm đoạt tài sản} kèm \textit{Giấy ủy quyền chuyển nhượng đất 200m2} mang tên Trần Ngọc Mai (ký hiệu \texttt{p2}).
\end{enumerate}

\section*{III. DANH MỤC VẬT CHỨNG THU THẬP NỘP KHO}
\begin{enumerate}[label=\arabic*., leftmargin=1.5cm]
    \item 01 Mảnh thủy tinh sắc nhọn dính máu khô (ký hiệu \texttt{p3}).
    \item 01 Cúc áo sơ mi nam màu xanh đen (ký hiệu \texttt{EV-SHIRT-BUTTON}).
    \item 01 Ghim cài áo huy hiệu công ty xây dựng (ký hiệu \texttt{p8}).
    \item Các mảnh giấy báo cũ năm 1998 bị xé rách (ký hiệu \texttt{p5}).
    \item Cuốn sổ ghi nợ bìa da đen bị xé trang (ký hiệu \texttt{10a}).
    \item Các tập hồ sơ giấy tờ tranh chấp đất 200m2, Đơn tố cáo của Mai (ký hiệu \texttt{p2}).
    \item 01 Mẫu máu khô và mẫu dấu giày nam size 41 thu tại khu vực bàn trà.
    \item 01 Hũ trà hoa cúc dán giấy note dặn uống trước khi ngủ.
\end{enumerate}

\noindent Biên bản khám nghiệm kết thúc vào hồi 11 giờ 30 phút cùng ngày, đã được đọc lại cho những người có tên trên nghe và cùng ký tên xác nhận.

\vspace{0.8cm}

% --- CHỮ KÝ NGHỊ ĐỊNH 30 ---
\noindent
\begin{minipage}[t]{0.45\textwidth}
    \small
    \textbf{\textit{Nơi nhận:}}\\
    - Thủ trưởng Cơ quan CSĐT (b/c);\\
    - VKSND TP. Hà Nội (để k/s);\\
    - Hồ sơ chuyên án TEST-99;\\
    - Lưu: VT, CSHS.
\end{minipage}
\hfill
\begin{minipage}[t]{0.5\textwidth}
    \begin{center}
        \small
        \textbf{CÁN BỘ CHỦ TRÌ KHÁM NGHIỆM}\\
        \textit{(Ký, ghi rõ họ tên)}\\[1.5cm]
        \textbf{Đại úy Lê Minh}
    \end{center}
\end{minipage}

\end{document}
